

\vspace{0.4cm}\emph{\textbf{(Lecture 5)}}\\

\begin{colordef}{Vector Field Along a Curve}{vecfield-along-curve}
Let $\gamma: I \to M$ be a smooth curve. We call a smooth map
\todo[inline]{ find a good notation for the bundle }
\[
X:t \longmapsto X_t \in T_{\gamma(t)} M
\]
a \emph{vector field along a curve} for $t \in I$, assigning to each $t$ a tangent vector at $\gamma(t)$.
\end{colordef}

\begin{colremark}
If $\gamma$ self-intersects (i.e., $\gamma(t_1) = \gamma(t_2)$), $X_{t_1}$ does not have to be the same as $X_{t_2}$ (so not a unique value at all points of $M$ belonging to the image of $\gamma$).

In this case: $X$ does not extend to a vector field on $M$. If $\gamma$ is an embedded curve, $X$ always extends (non-uniquely, i.e. in many different ways) to an element of $\Gamma(TM)$.

Locally (i.e., restricting the domain of $\gamma$): Any curve is embedded.
\end{colremark}



\begin{colorlem}{Covariant Derivative Along a Curve}{}
Let $\gamma: I \to M$ be a smooth curve and $X$ a vector field along $\gamma$. The covariant derivative $\nabla_{\dot{\gamma}} X$ is uniquely defined at each $t \in I$, independently of any local extension of $X$ around $\gamma(t)$.
\end{colorlem}

\begin{proof}[Sketch of proof]

Assuming $\gamma'(t)\neq 0$, restrict to $\gamma|_{(t-\epsilon,t+\epsilon)}$ so that it is embedded curve.



Extend $X$ to a vector field in a neighborhood of $\gamma(t)$ in $M$. Then, in local coordinates around $\gamma(t)$:
\begin{align*}
\left(\nabla_{\dot{\gamma}} X|_{\gamma(t)}\right)^k
&= \dot{\gamma}^i \frac{\partial}{\partial x^i} X^k + \Gamma_{ij}^k(\gamma(t))\ \dot{\gamma}^i X^j \\
&= \frac{\mathrm{d}}{\mathrm{d}t} X^i_t + \Gamma_{ij}^k(\gamma(t))\ \dot{\gamma}^i X^j_t
\end{align*}


The last expression depends only on the values along $\gamma$.
\end{proof}

\subsection*{Levi-Civita connection}

On a general manifold: An infinite family of connections. The difference of two such connections, $\nabla - \overline{\nabla}$ is always a (1,2)-tensor field (see exercise 4.1c). If $M$ is equipped with a Riemannian metric, there exists a unique \emph{natural} connection.

\begin{colortheo}{Levi-Civita Connection}{levi_civita}
Let $(M, g)$ be a Riemannian manifold. There exists a \textit{unique} connection $\nabla$ (called the Levi-Civita connection) such that:
\begin{enumerate}
    \item $\nabla$ is torsion-free: $\Gamma_{ij}^k = \Gamma_{ji}^k$ (if true in one coordinate system, it's true for all; since the torsion is a tensor).
    \item $\nabla$ respects the metric: for any $X, Y, Z \in \Gamma(TM)$,
    \[
    X(g(Y, Z)) = g(\nabla_X Y, Z) + g(Y, \nabla_X Z)
    \]
    (equivalently, $\nabla_X g = 0$ for all $X$.)
\end{enumerate}
\end{colortheo}

\begin{proof}


\textit{Uniqueness}: It suffices to prove that in any local coordinate system, conditions 1. \& 2. uniquely fix the Christoffel symbols.

In any local coordinates $(x^1, \ldots, x^n)$ we have

\begin{align*}
\frac{\partial}{\partial x^i} g_{jh} &= \frac{\partial}{\partial x^i}\left(g\left(\frac{\partial}{\partial x^j}, \frac{\partial}{\partial x^h}\right)\right) \\
&= g\left(\nabla_{\frac{\partial}{\partial x^i}} \frac{\partial}{\partial x^j}, \frac{\partial}{\partial x^h}\right) + g\left(\frac{\partial}{\partial x^j}, \nabla_{\frac{\partial}{\partial x^i}} \frac{\partial}{\partial x^h}\right) \\
&= g\left(\Gamma_{ij}^a \frac{\partial}{\partial x^a}, \frac{\partial}{\partial x^h}\right) + g\left(\frac{\partial}{\partial x^j}, \Gamma_{ih}^a \frac{\partial}{\partial x^a}\right) \\
& = \Gamma_{ij}^a g_{ah} + \Gamma_{ih}^a g_{aj}
\tag{4}
\end{align*}

Permuting indices:

\[
\frac{\partial}{\partial x^h} g_{ij} = \Gamma_{hj}^a g_{ai} + \Gamma_{hi}^a g_{aj} \quad \tag{1}
\]

\[
\frac{\partial}{\partial x^j} g_{ih} = \Gamma_{ji}^a g_{ah} + \Gamma_{jh}^a g_{ai} \quad \tag{2}
\]

So (1) + (2) - (4):

\[
\frac{\partial}{\partial x^i} g_{jh} + \frac{\partial}{\partial x^j} g_{ih} - \frac{\partial}{\partial x^h} g_{ij} = 2 \Gamma_{ij}^a g_{ah}
\]

So i can solve this for $\Gamma$, as the above is just a matrix product. We obtain, after renaming the indices,
\begin{equation}
    \label{eq:christoffel_levi_civita}
    \Gamma_{ij}^k = \frac{1}{2} g^{kd} \left( \frac{\partial}{\partial x^i} g_{jd} + \frac{\partial}{\partial x^j} g_{id} - \frac{\partial}{\partial x^d} g_{ij} \right)
\end{equation}

Where, as usual, $[g^{ij}]$ is the inverse matrix of $[g_{ij}]$. We see, that $\Gamma$ is uniquely determined by the metric $g$.\\

\textit{Existence}:
Instead of using the above formula to construct a connection, we go a slightly different path. Let us use the shorthand notation $\langle \cdot, \cdot \rangle := g(\cdot, \cdot)$. In order to define $\nabla_X Y$, $X, Y \in \Gamma(TM)$, it suffices to define $\langle \nabla_X Y, Z \rangle$, $\forall X, Y, Z \in \Gamma(TM)$.

(1)
\[
X \langle Y, Z \rangle = \langle \nabla_X Y, Z \rangle + \langle Y, \nabla_X Z \rangle
\]

(2)
\[
Y \langle X, Z \rangle = \langle \nabla_Y X, Z \rangle + \langle X, \nabla_Y Z \rangle
\]

(3)
\[
Z \langle X, Y \rangle = \langle \nabla_Z X, Y \rangle + \langle X, \nabla_Z Y \rangle
\]

We use now the torsion-free condition, $\nabla_X Y - \nabla_Y X = [X, Y]$.

From (1), (2), (3), we obtain Koszul's formula:

\[
\langle \nabla_X Y, Z \rangle
= \frac{1}{2} \Big(
    X(\langle Y, Z \rangle) + Y(\langle X, Z \rangle) - Z(\langle X, Y \rangle) + \langle [X, Y], Z \rangle + \langle [Z, X], Y \rangle - \langle [Y, Z], X \rangle \Big)
\]

One can use this formula to define a connection $\nabla$ on $M$. It is essentially a coordinate-free definition of, the Levi-Civita connection, while equation \eqref{eq:christoffel_levi_civita} is its coordinate expression.
\todo[inline]{ clean up the equation numbers }

\end{proof}


\begin{colexample}{Euclidean connection on $\mathbb{R}^n$ is Levi-Civita}{connection_Rn_is_levi_civita}
On $(\mathbb{R}^n, g_E)$, the standard directional derivative from equation \eqref{eq:connection_Rn}, $(D_X Y)^k = X^i \frac{\partial Y^k}{\partial x^i}$, is the Levi-Civita connection. The Christoffel symbols are zero in this case, as one can also verify using \eqref{eq:christoffel_levi_civita}.
\end{colexample}



\begin{colremark}[Levi-Civita Connection]
The Levi-Civita connection is the canonical connection on $(M, g)$. From now on, when referring to the connection of a Riemannian manifold, we will always mean the Levi-Civita connection.
\end{colremark}


Any construction based on a metric has to be preserved by isometries, as the latter identify geometric structures. In particular this holds for the Levi-Civita connection.
\begin{colorlem}{Isometries Preserve the Levi-Civita Connection}{}
If $F: (M, g) \longrightarrow (N, h)$ is a (local) isometry and $X, Y \in \Gamma(TM)$, then
\[
dF(\nabla_X^{(g)} Y) = \nabla_{dF(X)}^{(h)} dF(Y).
\]
\end{colorlem}

\begin{proof}

All the terms in the Koszul formula are "respected" by $dF$. I.e. we have $F$-isometry

\[
\langle dF(v), dF(w) \rangle_h = \langle v, w \rangle_g
\]

As well as
\[
[dF(v), dF(w)] = dF([v, w])
\]
for any smooth map.

So Koszuls formula stays the same under pushforwards by isometries, and therefore the Levi-Civita connection is preserved by isometries.
\todo[inline]{make this proof more readable}
\end{proof}




\begin{colorlem}{Christoffel symbols via orthogonal projection for an immersion}{}


Let $\Psi:\Omega\subset\mathbb{R}^n\to\mathbb{R}^N$ be a smooth immersion with induced metric $g=\Psi^* g_E$, and let $\Pi^\top$ denote orthogonal projection onto $T\Psi(\Omega)$. Then
\[
\Pi^\top\Big(\frac{\partial}{\partial x^i}\frac{\partial \Psi}{\partial x^j}\Big) = \Gamma^k_{ij}\,\frac{\partial\Psi}{\partial x^k}.
\]
\end{colorlem}

\begin{proof}
See exercise 5.4.
\end{proof}





\subsection{Parallel transport}

On a general manifold, parallel transport is the closest thing to a constant vector field along a curve.

\begin{colordef}{Parallel Transport}{}
    \todo[inline]{ make this definition of par. transp a bit more concise }
Let $\gamma: I \to M$ be smooth, and $\nabla$ be any connection. For $s,t\in I$, the \emph{parallel transport} along $\gamma$
\[
P_{s,t}:T_{\gamma(s)}M \longrightarrow T_{\gamma(t)}M,
\]
defined as follows: for $v\in T_{\gamma(s)}M$, let $X$ be the unique vector field along $\gamma$ such that
\[
\nabla_{\dot{\gamma}}X=0,\qquad X_s=v.
\]
Then
\[
P_{s,t}(v):=X_t.
\]
Equivalently, if $X$ is parallel along $\gamma$, then $X_t=P_{s,t}(X_s)$.
\end{colordef}

\begin{colortheo}{Existence and Uniqueness of Parallel Transport}{}
For every $p \in M$, $v \in T_pM$, and smooth curve $\gamma: I \to M$ with $\gamma(0) = p$, there exists a unique vector field $X$ along $\gamma$ such that $X_0 = v$ and $\nabla_{\dot{\gamma}} X = 0$.

\end{colortheo}

\begin{proof}
It suffices to do the construction in a local coordinate chart on $M$
\[
\frac{\mathrm{d}}{\mathrm{d}t} X^k + \Gamma_{ij}^k \dot{\gamma}^i X^j = 0
\]

This is a linear ODE, so ODE theory yields a unique solution.
\end{proof}

We note that the parallel transport has several nice properties, in particular for the Levi-Civita connection.
\begin{colorlem}{Basic Properties of Parallel Transport}{basic_properties_parallel_transport}
Let $\gamma:I\to M$ be smooth, and let $P_{s,t}:T_{\gamma(s)}M\longrightarrow T_{\gamma(t)}M$ denote parallel transport along $\gamma$.

\begin{enumerate}
    \item For all $s,t\in I$, $P_{s,t}$ is a linear isomorphism, with inverse $P_{t,s}$.
    \item If $\nabla$ is the Levi-Civita connection of $(M,g)$, then $P_{s,t}$ is an isometry:
    \[
    g_{\gamma(t)}\!\big(P_{s,t}v,P_{s,t}w\big)=g_{\gamma(s)}(v,w),
    \qquad v,w\in T_{\gamma(s)}M.
    \]
\end{enumerate}
\end{colorlem}

\begin{proof}
Parallel transport along $\gamma$ is defined by a linear ODE in each coordinate chart,
\[
\frac{\mathrm{d}}{\mathrm{d}t} X^k + \Gamma_{ij}^k(\gamma(t))\, \dot{\gamma}^i X^j = 0.
\]
The equation is linear in $X$, so the solution depends linearly on the initial condition
$v \in T_{\gamma(0)}M$. By existence and uniqueness for linear ODEs, for each $v$ there is a unique
vector field $X(t)$ along $\gamma$.

Hence the map
\[
P_t : v \longmapsto X(t)
\]
is linear and bijective, with inverse given by solving the ODE backward in time.\\

Concerning the isometry property, let $X, Y$ be vector fields on $\gamma$ s.t. $\nabla_\gamma X = \nabla_\gamma Y = 0$. Then
$$\frac{\mathrm{d}}{\mathrm{d}t} \left \langle X,Y \right \rangle = \gamma' (\left \langle X,Y \right \rangle) = \left \langle \nabla_\gamma X,Y \right \rangle + \left \langle X,\nabla_\gamma Y \right \rangle = 0,$$

where we have used that for $f(p):=\left \langle X_p,Y_p \right \rangle$ we have $\frac{\mathrm{d}}{\mathrm{d}t}(f\circ \gamma)(t) = \mathrm{d}\gamma \left ( \frac{\mathrm{d}}{\mathrm{d}t} \right )(f)=\gamma'(t)f$. 

The map $t\mapsto \left \langle X(t),Y(t) \right \rangle$ lives on $I\subset \mathbb{R}$, and has zero derivative, so we get $\left \langle X(t),Y(t) \right \rangle = \left \langle X(0),Y(0) \right \rangle$, which is exactly the isometric property.
\end{proof}

In particular, if $\gamma(a)=\gamma(b)$ then $P_t:T_{\gamma(a)} M \to T_{\gamma(b)} M$ is the identity map, so parallel transport along a closed curve is an isometry of the tangent space at the base point.




















\newpage
\section{Geodesics and the exponential map}

\subsection{Geodesics}

We define geodesics as curves whose velocity is parallel along the curve.

\begin{colordef}{Geodesic}{geodesic}
A geodesic is a curve $\gamma: I \to M$ that satisfies
\[
\nabla_{\dot{\gamma}} \dot{\gamma} = 0.
\]
\end{colordef}

In local coordinates, this yields the \emph{geodesic equations},
\begin{equation}
    \label{eq:geodesic_equation}
    \ddot{\gamma}^k + \Gamma_{ij}^k \dot{\gamma}^i \dot{\gamma}^j = 0
\end{equation}

where $\Gamma_{ij}^k$ are the Christoffel symbols of the connection.

\begin{colremark}
If we reparametrize: $\gamma' = \gamma \circ h$: one can show that $\nabla_{\dot{\gamma}'} \dot{\gamma}' \parallel \dot{\gamma}'$, i.e. there exists a function $\lambda(t)$ such that $\nabla_{\dot{\gamma}'} \dot{\gamma}' =\lambda \dot{\gamma}'$. But we do not necessarily have $\nabla_{\dot{\gamma}'} \dot{\gamma}'=0$ anymore so the above definition is parametrisation dependent. Any curve which satisfies $\nabla_{\dot{\gamma}'} \dot{\gamma}' =\lambda \dot{\gamma}'$ can always be reparametrised to geodesics.
\end{colremark}



Consider now an embedded hypersurface $S$ in $\mathbb{R}^n$ with the induced metric, and the Levi-Civita connection $\nabla$ of $g$. Recall further example \ref{ex:connection_Rn_is_levi_civita}, and \eqref{eq:connection_Rn}. One can show that on $S\subset \mathbb{R}^n$, the Levi-Civita connection is the projection of the Euclidean connection on $\mathbb{R}^n$ onto $TS$ (to do so, define $\nabla_X Y:=\Pi(D_X Y)$ and verify the assumptions of Theorem \ref{thm:levi_civita}). So we have
\[\nabla_{\dot{\gamma}} \dot{\gamma} = \Pi(D_{\dot{\gamma}} \dot{\gamma}),\]
where $\Pi$ is the $g$-orthogonal projection onto the tangent space $T_{\gamma(t)}S$, and where $D$ is Euclidean connection on $\mathbb{R}^n$ as in equation \eqref{eq:connection_Rn}.

In this setting, a curve $\gamma:(a,b) \to S$ is a geodesic iff
\[
0 = \nabla_{\dot{\gamma}} \dot{\gamma} = \Pi(D_{\dot{\gamma}} \dot{\gamma}),
\]
where $\ddot\gamma:=D_{\dot{\gamma}} \dot{\gamma} = \frac{\mathrm{d}^2\gamma}{\mathrm{d}t^2}$. So here, the \emph{ambient acceleration in} $\mathbb{R}^n$ of a geodesic is orthogonal to the surface, $\ddot\gamma\perp T_{\gamma(t)}S$.


\begin{colexample}{Geodesics on \((S^n, g_{S^n})\)}{}
Consider \((S^n, g_{S^n})\), where \(g_{S^n}\) is the metric induced from \((\mathbb{R}^{n+1}, g_e)\). In view of the previous discussion, geodesics \(\gamma\) are exactly curves whose acceleration in \(\mathbb{R}^{n+1}\) is orthogonal to \(TS^n\). Since \(T_{\gamma(t)}S^n = \gamma(t)^\perp\), the condition \(\ddot\gamma \perp T_{\gamma(t)}S^n\) implies \(\ddot\gamma = \lambda(t)\gamma\). Hence the subspace
\[
V := \mathrm{span}\{\gamma(t), \dot\gamma(t)\}
\]
is invariant under time differentiation, so \(\gamma(t) \subset V\) for all \(t\). Therefore \(\gamma\) lies in a 2-dimensional linear subspace through the origin, and thus is a great circle.
\end{colexample}


\begin{center}
\begin{tikzpicture}[scale=2]
    % Draw circle (S^1)
    \draw[thick] circle (1);
    
    % Mark point p on circle
    \draw[fill=black] (0,1) circle (0.05);
    \node[above] at (0,1) {$\gamma(t)$};
    
    % Draw tangent line at p
    \draw[thick, Groseille] (-2,1) -- (2,1);
    
    % Draw tangent vector
    \draw[->, thick, Canard, shift={(0,1)}] (0,0) -- (1,0);
    \node[right, above, Canard] at (1,1) {$\dot\gamma$};

    \node[right, above, Groseille] at (-1,1) {$T_{\gamma(t)} S^{n}$};
    \draw[fill=black] (0,0) circle (0.05);
    \node[right, below, black] at (0,0) {$0\in\mathbb{R}^{n+1}$};
    
    % acceleration
    \draw[->, thick, Canard, shift={(0,0)}] (0,1) -- (0,0.5);
    \node[right, Canard] at (0,0.5) {$\ddot\gamma$};

    % Image of exp_p(v) on sphere (projected along circle)
    % \coordinate (q) at (0.8,0.6); % on circle
    % \fill (q) circle (0.8pt) node[right] {$\exp_p(v)$};
    
    % Label the circle
    \node at (1,-0.7) {$S^1$};
    
    % Draw normal vector (radius)
    \draw[dashed, gray] (0,0) -- (0,1);
\end{tikzpicture}
\end{center}

\begin{colortheo}{Existence and Uniqueness of Geodesics}{existence_uniqueness_geodesics}

For every $p \in M$ and $v \in T_pM$, there exists a unique maximal (domain can not be extended) geodesic $\gamma: (-\epsilon, \epsilon) \to M$ such that $\gamma(0) = p$ and $\dot{\gamma}(0) = v$. 
\end{colortheo}

For such $p,v$ we will denote the curve by $\gamma_{p,v}(t)$ from now on.

\begin{proof}[Proof sketch]
ODE theory yields existence and uniqueness of solutions to the geodesic equations, \eqref{eq:geodesic_equation}, and that $\gamma$ depends smoothly on $p$ and $v$. 
\end{proof}

\begin{colorcor}{}{}
If you have two curves $\gamma_1, \gamma_2$ such that $\gamma_1(0) = \gamma_2(0)$ and $\dot{\gamma}_1(0) = \dot{\gamma}_2(0)$, then on the common domain of intersection they have to agree, i.e. they are subcurves of the same maximal geodesic.
\end{colorcor}



\begin{colremark}
\label{rem:euclidean_geodesics_are_straight}
Recall that $\Gamma_{ij}^k=0$ in cartesian coordinates due to \eqref{eq:christoffel_levi_civita}. Hence geodesics satisfy $\ddot{\gamma}^k + \Gamma_{ij}^k \dot{\gamma}^i \dot{\gamma}^j = \ddot{\gamma}^k = 0$, so by integrating twice we get $\gamma^k(t) = a^k t + b^k$, i.e. geodesics are straight lines in $\mathbb{R}^n$. Hence they minimize the distance between points.

On a general Riemannian manifold, this is not necessarily true (for example, on the sphere). However, geodesics are always stationary points of the length functional, as we will see later.
\todo[inline]{ link stationary points and local minimsaion }
\end{colremark}

From the definition, we can easily show that $\left\lVert \gamma' \right\rVert=g(\gamma',\gamma')$ is constant in $t$, i.e. geodesics are parametrised at constant speed.

\begin{colorlem}{Geodesics are parametrised at constant speed}{geodesics_constant_speed}
Let $\gamma$ be a geodesic on $(M,g)$, and let $\nabla$ be the Levi-Civita connection. Then 
$$\frac{\mathrm{d}}{\mathrm{d}t} g(\gamma',\gamma')=0.$$
\end{colorlem}
\begin{proof}
    \begin{align*}
        \frac{\mathrm{d}}{\mathrm{d}t} g(\gamma',\gamma')
        &= \mathrm{d} \gamma \left ( \frac{\mathrm{d}}{\mathrm{d}t} \right ) \left ( g(\gamma',\gamma') \right )\\
        &= \gamma'(t) g(\gamma',\gamma')\\
        &= g(\nabla_{\gamma'}\gamma',\gamma') + g(\gamma',\nabla_{\gamma'}\gamma'). 
    \end{align*}
    But $\nabla_{\gamma'}\gamma'=0$ as $\gamma$ is a geodesic, so we yield the result.
\end{proof}



We present now a version of Nöther's Theorem.

\begin{colorlem}{Nöther's Theorem for Geodesics}{noether}
Let $X \in \Gamma(TM)$ be Killing, $\nabla$ the Levi-Civity connection, and $\gamma$ be a Geodesic on $(a,b)$. Then the function $t \mapsto g(X|_{\gamma(t)}, \dot{\gamma})$ is constant for $t \in (a, b)$.
\end{colorlem}

\begin{proof}
See exercise 5.1a.
\end{proof}

Any function $F: TM \to \mathbb{R}$ such that $F(\gamma(t), \dot{\gamma}(t))$ is constant when $\gamma$ is a geodesic is called a \emph{constant of motion} for the geodesic flow. Classically, Nöther's theorem states that for any \emph{symmetry} of the system, there is a corresponding constant of motion. In this case, the function $t \mapsto g(X|_{\gamma(t)}, \dot{\gamma})$ is a constant of motion along the geodesic. The symmetry is given by the Killing vector field $X$, which, as we have seen is related to isometries of the manifold, see Lemma \ref{lem:killing_generates_isometries}.

We will see later, that isometries map geodesics to geodesics.



\begin{colexample}{Geodesics and constants of motion in the Poincaré Half-Plane}{}
Let $\mathbb{H}^2:=\{(x,y):y>0\}$ be the Poincaré half-plane with $g_H = \frac{dx^2 + dy^2}{y^2}$. In exercise 5.2 we show that 
\begin{itemize}[itemsep=0pt, topsep=0pt]
    \item Maps $f(z) = \frac{az+b}{cz+d}$, with $ad-bc>0$, are isometries
    \item The geodesic equations are given by 
    \[
    \ddot{x} - \frac{2}{y} \dot{x}\dot{y} = 0, \quad 
    \ddot{y} + \frac{1}{y} (\dot{x}^2 - \dot{y}^2) = 0,
    \]
    with conserved quantities
    \[
    E = g_H(\dot{\gamma}, \dot{\gamma}), \quad P = \frac{\dot{x}}{y^2},
    \]  
    P being a constant of motion associated to the Killing vector field $\frac{\partial}{\partial x}$, as in Lemma \ref{lem:noether}.
\item The geodesics are semicircles orthogonal to the x-axis.
\end{itemize}

\begin{center}
\begin{tikzpicture}[scale=1.0,>=stealth]

    % Axes and boundary
    \draw[<->] (-4.2,0) -- (4.3,0) node[right] {$x$};
    \draw[->] (0,0) -- (0,3.35) node[above] {$y$};

    % Geodesics in H^2: vertical lines and semicircles orthogonal to y=0
    \draw[->,very thick,Rouge] (0.0,0) -- (0.0,3.0);
    \draw[<-,very thick,Groseille] (-3.2,0) arc[start angle=180,end angle=0,radius=1.6];
    \draw[->,very thick,Canard] (0.0,0) arc[start angle=180,end angle=0,radius=2.0];
    \draw[->,very thick,Leman] (0,0) arc[start angle=180,end angle=0,radius=0.9];

    \node at (3,3) {\small $\mathbb H^2=\{(x,y):y>0\}$};
\end{tikzpicture}
\end{center}

\end{colexample}



\begin{colorlem}{Parallel orthonormal frame along a geodesic}{parallel_orthonormal_frame}
Let $\gamma:[0,1]\to M$ be a geodesic. Then there exists an orthonormal frame $\{E_i\}_{i=1}^n$ along $\gamma$ such that
\[
E_1(t)\ \text{is collinear with }\dot{\gamma}(t), 
\qquad
\nabla_{\dot{\gamma}}E_i=0 \quad \text{for all } i=1,\dots,n.
\]

Moreover, for any vector field $X=X_i(t)\,E_i|_{\gamma(t)}$ parallel-transported along $\gamma$, one has
\[
\nabla_{\dot{\gamma}}X=0
\quad\Longleftrightarrow\quad
X_i(t) \text{ is constant in $t$ for all } i.
\]
\end{colorlem}

\begin{proof}
See exercise 5.3.
\end{proof}



